\documentclass{article}
\usepackage{amsmath}
% \documentclass{article}
\usepackage{amsmath}

\begin{document}

\textbf{Ejercicio 41}

\[
\int_{-2}^{2} \frac{dx}{\sqrt{4 - x^2}}
\]

\textbf{Solución}

Sabemos que:
\[
\int \frac{dx}{\sqrt{a^2 - x^2}} = \arcsin\left(\frac{x}{a}\right) + C
\]

En este caso:
\[
a^2 = 4 \Rightarrow a = 2
\]

Entonces:
\[
\int \frac{dx}{\sqrt{4 - x^2}}
= \arcsin\left(\frac{x}{2}\right)
\]

Evaluando en los límites:
\[
\left[ \arcsin\left(\frac{x}{2}\right) \right]_{-2}^{2}
= \arcsin(1) - \arcsin(-1)
= \frac{\pi}{2} + \frac{\pi}{2}
= \pi
\]

\[
\boxed{\pi}
\]


\textbf{Ejercicio 42}

\[
\int_{0}^{4} \frac{dx}{4 - x}
\]

\textbf{Solución}

Observamos que la función no está definida en \(x = 4\), por lo tanto
la integral es impropia.

\[
\int_{0}^{4} \frac{dx}{4 - x}
= \lim_{b \to 4^-} \int_{0}^{b} \frac{dx}{4 - x}
\]

Calculamos la integral:
\[
\int \frac{dx}{4 - x} = -\ln|4 - x| + C
\]

Evaluamos en los límites:
\[
\lim_{b \to 4^-}
\left[ -\ln|4 - x| \right]_{0}^{b}
\]

\[
= \lim_{b \to 4^-}
\left( -\ln(4 - b) + \ln 4 \right)
\]

Como \(4 - b \to 0^+\), se tiene:
\[
\ln(4 - b) \to -\infty
\]

Por lo tanto:
\[
-\ln(4 - b) \to +\infty
\]

---

\[
\boxed{\int_{0}^{4} \frac{dx}{4 - x} \text{ diverge}}
\]

\textbf{Ejercicio 43}

\[
\int_{-\infty}^{0} x e^x \, dx
\]

\textbf{Solución}

La integral es impropia porque el límite inferior es \(-\infty\).

\[
\int_{-\infty}^{0} x e^x \, dx
= \lim_{a \to -\infty} \int_{a}^{0} x e^x \, dx
\]

Usamos integración por partes:

\[
\begin{aligned}
u &= x \quad & dv &= e^x dx \\
du &= dx \quad & v &= e^x
\end{aligned}
\]

Entonces:
\[
\int x e^x \, dx = x e^x - \int e^x dx
\]

\[
= x e^x - e^x
\]

\[
= e^x (x - 1)
\]

Evaluamos en los límites:
\[
\lim_{a \to -\infty}
\left[ e^x (x - 1) \right]_{a}^{0}
\]

\[
= \left( e^0 (0 - 1) \right)
- \lim_{a \to -\infty} e^a (a - 1)
\]

\[
= -1 - 0
\]

---

\[
\boxed{\int_{-\infty}^{0} x e^x \, dx = -1}
\]


\textbf{Ejercicio 44}

\[
\int_{-\infty}^{+\infty} \frac{dx}{1 + 4x^2}
\]

\textbf{Solución}

La integral es impropia porque los límites son infinitos.

\[
\int_{-\infty}^{+\infty} \frac{dx}{1 + 4x^2}
= \lim_{a \to -\infty} \int_{a}^{0} \frac{dx}{1 + 4x^2}
+ \lim_{b \to +\infty} \int_{0}^{b} \frac{dx}{1 + 4x^2}
\]

Sabemos que:
\[
\int \frac{dx}{a^2 + x^2} = \frac{1}{a} \arctan\left(\frac{x}{a}\right) + C
\]

Reescribimos:
\[
\frac{1}{1 + 4x^2} = \frac{1}{1^2 + (2x)^2}
\]

Entonces:
\[
\int \frac{dx}{1 + 4x^2}
= \frac{1}{2} \arctan(2x)
\]

Evaluamos los límites:
\[
\left[ \frac{1}{2} \arctan(2x) \right]_{-\infty}^{+\infty}
\]

\[
= \frac{1}{2}
\left( \frac{\pi}{2} - \left(-\frac{\pi}{2}\right) \right)
\]

\[
= \frac{1}{2} \pi
\]

---

\[
\boxed{\int_{-\infty}^{+\infty} \frac{dx}{1 + 4x^2} = \frac{\pi}{2}}
\]



% \begin{document}

\textbf{Ejercicio 45}

\[
y^2 = 8x^2 \quad \text{desde } x=1 \text{ hasta } x=8
\]

Despejamos \(y\):
\[
y = 2\sqrt{2}\,x
\]

La longitud de arco es:
\[
L=\int_a^b \sqrt{1+\left(\frac{dy}{dx}\right)^2}\,dx
\]

Derivando:
\[
\frac{dy}{dx} = 2\sqrt{2}
\]

Sustituyendo:
\[
L=\int_{1}^{8} \sqrt{1+(2\sqrt{2})^2}\,dx
= \int_{1}^{8} \sqrt{9}\,dx
\]

\[
L=\int_{1}^{8} 3\,dx
=3(8-1)
=21
\]

\[
\boxed{L = 21}
\]

\textbf{Ejercicio 46}

\[
27y^2 = 4(x-2)^3
\quad \text{desde } (2,0) \text{ hasta } (11,6\sqrt{3})
\]

Despejamos:
\[
y = \frac{2}{3\sqrt{3}}(x-2)^{3/2}
\]

Derivamos:
\[
\frac{dy}{dx} = \frac{1}{\sqrt{3}}(x-2)^{1/2}
\]

Longitud de arco:
\[
L=\int_{2}^{11} \sqrt{1+\left(\frac{dy}{dx}\right)^2}\,dx
\]

\[
= \int_{2}^{11} \sqrt{\frac{x+1}{3}}\,dx
= \frac{1}{\sqrt{3}}\int_{2}^{11} \sqrt{x+1}\,dx
\]

\[
= \frac{2}{3\sqrt{3}}
\left[(x+1)^{3/2}\right]_{2}^{11}
\]

\[
= \frac{2}{3\sqrt{3}}(24\sqrt{3}-3\sqrt{3})
= 14
\]

\[
\boxed{L=14}
\]



\textbf{Ejercicio 47}

\[
y=\ln x \quad \text{desde } x=1 \text{ hasta } x=2\sqrt{2}
\]

Derivamos:
\[
\frac{dy}{dx}=\frac{1}{x}
\]

Longitud de arco:
\[
L=\int_{1}^{2\sqrt{2}} \sqrt{1+\frac{1}{x^2}}\,dx
=\int_{1}^{2\sqrt{2}} \frac{\sqrt{x^2+1}}{x}\,dx
\]

La primitiva es:
\[
\sqrt{x^2+1}
+\ln\!\left(\frac{x}{\sqrt{x^2+1}+1}\right)
\]

Evaluando:
\[
L=\left[
\sqrt{x^2+1}
+\ln\!\left(\frac{x}{\sqrt{x^2+1}+1}\right)
\right]_{1}^{2\sqrt{2}}
\]

\[
\boxed{
L=3-\sqrt{2}
+\ln\!\left(\frac{(2\sqrt{2})(\sqrt{2}+1)}{4}\right)
}
\]

\textbf{Ejercicio 48}

\[
y=\ln\!\left(\frac{e^x-1}{e^x+1}\right),
\qquad 2\le x\le 4
\]

Derivamos:
\[
\frac{dy}{dx}
=\frac{e^x}{e^x-1}-\frac{e^x}{e^x+1}
=\frac{2e^x}{e^{2x}-1}
\]

Longitud de arco:
\[
L=\int_{2}^{4}\sqrt{1+\left(\frac{dy}{dx}\right)^2}\,dx
\]

\[
\sqrt{1+\left(\frac{2e^x}{e^{2x}-1}\right)^2}
=\frac{e^{2x}+1}{e^{2x}-1}
\]

\[
L=\int_{2}^{4}\frac{e^{2x}+1}{e^{2x}-1}\,dx
\]

\[
= \int_{2}^{4}dx
+2\int_{2}^{4}\frac{dx}{e^{2x}-1}
\]

\[
L=\left[\ln(e^{2x}-1)-x\right]_{2}^{4}
\]

\[
\boxed{
L=\ln\!\left(\frac{e^8-1}{e^4-1}\right)-2
}
\]

\textbf{Ejercicio 49}

\[
y=\ln(\cos x), \qquad \frac{\pi}{6}\le x\le \frac{\pi}{4}
\]

Derivamos:
\[
\frac{dy}{dx}=-\tan x
\]

Longitud de arco:
\[
L=\int_{\pi/6}^{\pi/4}\sqrt{1+\tan^2 x}\,dx
\]

Usando \(1+\tan^2 x=\sec^2 x\):
\[
L=\int_{\pi/6}^{\pi/4}\sec x\,dx
\]

\[
L=\left[\ln(\sec x+\tan x)\right]_{\pi/6}^{\pi/4}
\]

\[
\boxed{
L=\ln\!\left(\frac{\sqrt{2}+1}{\sqrt{3}}\right)
}
\]


\textbf{Ejercicio 50}

Calcular la longitud del arco de una circunferencia de radio \(3\)
que forma un ángulo de \(60^\circ\) con el eje \(x\) positivo.

\bigskip

\textbf{Método 1: Longitud de arco}

La circunferencia se parametriza como:
\[
x=3\cos t, \qquad y=3\sin t
\]

con:
\[
t\in\left[0,\frac{\pi}{3}\right]
\]

La longitud de arco es:
\[
L=\int_0^{\pi/3}
\sqrt{\left(\frac{dx}{dt}\right)^2+\left(\frac{dy}{dt}\right)^2}\,dt
\]

\[
= \int_0^{\pi/3} 3\,dt
=3\left[\frac{\pi}{3}\right]
=\pi
\]

\bigskip

\textbf{Método 2: Fórmula del arco de circunferencia}

\[
L=r\theta
\]

\[
\theta=60^\circ=\frac{\pi}{3}
\]

\[
L=3\cdot\frac{\pi}{3}=\pi
\]

\[
\boxed{L=\pi}
\]

\end{document}